\documentclass[11pt,a4paper]{article}

\usepackage[margin=1in]{geometry}
\usepackage{hyperref}
\usepackage{listings}
\usepackage{xcolor}
\usepackage{enumitem}

\definecolor{codegray}{rgb}{0.95,0.95,0.95}

\lstset{
    backgroundcolor=\color{codegray},
    basicstyle=\ttfamily\small,
    frame=single,
    breaklines=true,
    showstringspaces=false
}
% --- Custom languages for listings (TypeScript + Prisma) ---
\lstdefinelanguage{TypeScript}{
    keywords={break, case, catch, class, const, continue, debugger, default, delete, do, else, enum, export, extends, false, finally, for, function, if, import, in, instanceof, new, null, return, super, switch, this, throw, true, try, typeof, var, void, while, with, as, implements, interface, let, package, private, protected, public, static, yield, any, boolean, number, string, symbol, type, from, of},
    sensitive=true,
    morecomment=[l]{//},
    morecomment=[s]{/*}{*/},
    morestring=[b]',
    morestring=[b]"
}

\lstdefinelanguage{Prisma}{
    keywords={generator, datasource, model, enum, type, provider, url, env, relation, map, db, String, Int, Boolean, DateTime, Float, Decimal, Json, Bytes},
    sensitive=true,
    morecomment=[l]{//},
    morecomment=[s]{/*}{*/},
    morestring=[b]"
}


\title{\textbf{FlexInvoice Project Flow}}
\author{Backend Development Log}
\date{\today}

\begin{document}
    \maketitle
    \tableofcontents
    \newpage

% --------------------------------------------------
    \section{Project Overview}

    FlexInvoice is a backend system for invoice, sales, and inventory management.
    The backend is implemented using a RESTful architecture with a strong focus on:

    \begin{itemize}
        \item Security-by-design
        \item Clean, maintainable structure
        \item Incremental development
        \item Database-first correctness
    \end{itemize}

    The project is intentionally built in clearly defined steps to ensure correctness
    before feature complexity.

% --------------------------------------------------
    \section{Step 00 -- Foundations}

    \subsection{Objective}

    The goal of this step was to establish a minimal but production-grade backend
    foundation before implementing any business logic.

    Specifically:
    \begin{itemize}
        \item Ensure the server boots reliably
        \item Validate environment configuration at runtime
        \item Establish secure middleware defaults
        \item Avoid premature abstractions or clutter
    \end{itemize}

    \subsection{Technology Stack}

    \begin{itemize}
        \item Node.js (LTS)
        \item TypeScript
        \item Express.js (REST API)
        \item Zod (runtime validation)
        \item Pino (structured logging)
    \end{itemize}

    \subsection{Key Design Decisions}

    \begin{itemize}
        \item Express was chosen over opinionated frameworks to retain control.
        \item Environment variables are validated at startup; the app will not run with invalid config.
        \item No controllers/services/repositories folders were created prematurely.
    \end{itemize}

    \subsection{Implementation Summary}

    \begin{itemize}
        \item \texttt{src/app.ts} initializes middleware in a strict order:
        logging $\rightarrow$ security headers $\rightarrow$ CORS $\rightarrow$ parsing $\rightarrow$ routes.
        \item \texttt{src/server.ts} is responsible only for starting the HTTP server.
        \item A health endpoint (\texttt{/health}) was added for diagnostics and deployment readiness.
    \end{itemize}

    \subsection{Security Measures Applied}

    \begin{itemize}
        \item HTTP security headers via Helmet
        \item Request body size limits
        \item Sensitive fields redacted from logs
        \item Centralized error handler preventing stack trace leaks
    \end{itemize}

% --------------------------------------------------
    \section{Step 01 -- Prisma Schema and MySQL Integration}

    \subsection{Objective}

    The objective of this step was to integrate a relational database layer
    in a way that is:
    \begin{itemize}
        \item Type-safe
        \item Migration-driven
        \item Easy to reason about
        \item Replaceable later (MySQL $\rightarrow$ MongoDB)
    \end{itemize}

    \subsection{Database Choice}

    MySQL was selected for the initial phase due to:
    \begin{itemize}
        \item Strong relational guarantees
        \item Developer familiarity
        \item Mature tooling
    \end{itemize}

    A future migration to MongoDB is planned once the domain stabilizes.

    \subsection{Prisma Configuration}

    Prisma was configured using \texttt{prisma.config.ts}, which explicitly loads
    environment variables using \texttt{dotenv}. This avoids implicit CLI behavior
    and ensures deterministic configuration.

    \begin{lstlisting}[language=TypeScript]
import 'dotenv/config';
import { defineConfig, env } from 'prisma/config';

export default defineConfig({
  schema: 'prisma/schema.prisma',
  datasource: {
    url: env('DATABASE_URL'),
  },
});
    \end{lstlisting}

    \subsection{Schema Design}

    The initial schema introduced a \texttt{User} model and a \texttt{Role} enum:

    \begin{lstlisting}[language=Prisma]
model User {
  id        Int      @id @default(autoincrement())
  email     String   @unique
  password  String
  role      Role     @default(USER)
  createdAt DateTime @default(now())
  updatedAt DateTime @updatedAt
}

enum Role {
  USER
  SELLER
  ADMIN
}
    \end{lstlisting}

    \subsection{Migration Strategy}

    \begin{itemize}
        \item Prisma migrations are committed to version control.
        \item Generated Prisma Client code is ignored.
        \item Database changes are never applied manually.
    \end{itemize}

    The initial migration was successfully applied using:
    \begin{lstlisting}
npx prisma migrate dev --name init
    \end{lstlisting}

    \subsection{Verification}

    A temporary debug endpoint was added to verify end-to-end connectivity:

    \begin{itemize}
        \item Express routing
        \item Prisma Client generation
        \item MySQL connectivity
    \end{itemize}

    The endpoint returned an empty user list, confirming correct setup.

    \subsection{Lessons Learned}

    \begin{itemize}
        \item Prisma CLI and application runtime load environment variables independently.
        \item Passwords containing special characters must be URL-encoded in database URLs.
        \item Node.js bleeding-edge versions may cause compatibility issues with ORM tooling.
    \end{itemize}

% --------------------------------------------------
    \section{Current State Summary}

    At the end of this phase:
    \begin{itemize}
        \item The backend boots reliably
        \item Environment configuration is validated
        \item Prisma is fully operational with MySQL
        \item Migrations are tracked and reproducible
        \item No business logic has been implemented yet
    \end{itemize}

    The project is now ready to implement authentication and authorization.

\end{document}
